\chapter{The Setup on This Machine}\label{ch:vim-setup}

The previous chapter described the pieces in the abstract. This one describes
the Neovim installation these notes were written in, so that the abstractions
have somewhere concrete to land. At the time of writing that is Neovim 0.9.5
running \vocab{LazyVim} with VimTeX and LuaSnip added on top.

\section{LazyVim}\label{sec:setup-lazyvim}

\begin{definition}[Distribution]
	A pre-assembled configuration: a plugin manager, a curated plugin set, and a
	consistent set of keybindings, packaged so it can be installed in one step and
	then overridden piecemeal. \texttt{LazyVim} \cite{lazyvim} is the most widely used one.
\end{definition}

The layout it expects is small:

\begin{vimcode}
~/.config/nvim/
  init.lua                 # entry point; requires config.lazy
  lazy-lock.json           # pinned plugin commits
  lua/config/lazy.lua      # bootstraps lazy.nvim and imports LazyVim
  lua/config/options.lua   # your option overrides
  lua/config/keymaps.lua   # your mappings
  lua/config/autocmds.lua  # your autocommands
  lua/plugins/*.lua        # one file per plugin or topic; merged automatically
\end{vimcode}

Every file under \file{lua/plugins/} is imported and its specs are merged with
LazyVim's. Declaring the same plugin again \emph{extends} the existing spec
rather than replacing it, which is how you change one option without forking
the whole configuration:

\begin{luacode}
-- lua/plugins/ui.lua: adjust a plugin LazyVim already installs
return {
  {
    "folke/tokyonight.nvim",
    opts = { transparent = true },
  },
}
\end{luacode}

\begin{note}[Discovering the keymap]
	LazyVim binds several hundred keys and you are not meant to memorise them.
	Press \key{<Space>} and stop --- \texttt{which-key} pops up a menu of every
	continuation, grouped and labelled. \key{<leader>} \key{s} \key{k} searches
	all mappings, and \cmd{:verbose nmap <keys>} says which file defined one.
\end{note}

\section{The Keymap Worth Learning First}\label{sec:setup-keys}

Verified against the LazyVim installed here; the leader is \key{<Space>}.

\begin{keys}[0.26]
	\key{<leader>} \key{<Space>} & Find files in the project root        \\
	\key{<leader>} \key{/}       & Grep across the project               \\
	\key{<leader>} \key{e}       & File explorer                         \\
	\key{<leader>} \key{,}       & Switch buffer                         \\
	\key{<S-h>} / \key{<S-l>}    & Previous / next buffer                \\
	\key{<leader>} \key{`}       & Toggle with the alternate buffer      \\
	\key{<leader>} \key{b} \key{d} & Delete this buffer, keep the window \\
	\key{<C-h>} \key{<C-j>} \key{<C-k>} \key{<C-l>} & Move between windows (no \key{<C-w>} prefix) \\
	\key{<leader>} \key{-} / \key{<leader>} \key{|} & Split below / right \\
	\key{<C-/>}                  & Toggle a floating terminal            \\
	\key{<C-s>}                  & Save, from any mode                   \\
	\key{<leader>} \key{q} \key{q} & Quit everything                     \\
	\key{<leader>} \key{q} \key{s} & Restore this directory's session    \\
	\key{<leader>} \key{l}       & Open the \cmd{:Lazy} dashboard        \\
	\key{<leader>} \key{u} \key{i} & Inspect highlight groups under the cursor \\
\end{keys}

Code-related, active whenever a language server is attached:

\begin{keys}[0.26]
	\key{K}                        & Hover documentation             \\
	\key{g} \key{d} / \key{g} \key{D} & Definition / declaration      \\
	\key{g} \key{r}                & References                      \\
	\key{g} \key{I} / \key{g} \key{y} & Implementation / type definition \\
	\key{g} \key{K}                & Signature help                  \\
	\key{<leader>} \key{c} \key{a} & Code action                     \\
	\key{<leader>} \key{c} \key{r} & Rename symbol                   \\
	\key{<leader>} \key{c} \key{f} & Format buffer or selection      \\
	\key{<leader>} \key{c} \key{d} & Line diagnostics in a float     \\
	\key{]} \key{d} / \key{[} \key{d} & Next / previous diagnostic   \\
	\key{]} \key{e} / \key{[} \key{e} & Next / previous \emph{error} \\
	\key{<leader>} \key{x} \key{x} & Diagnostics list (Trouble)      \\
	\key{<leader>} \key{s} \key{r} & Project-wide search and replace (grug-far) \\
	\key{g} \key{c} \key{c} / \key{g} \key{c} & Comment a line / a motion \\
	\key{<A-j>} / \key{<A-k>}      & Move the current line or selection down / up \\
	\key{s}                        & Flash jump: type two characters, then a label \\
\end{keys}

\begin{note}
	LazyVim also remaps \key{j} and \key{k} to \key{g} \key{j} and \key{g} \key{k}
	when no count is given, so cursor movement follows display lines in wrapped
	prose but \key{5} \key{j} still means five real lines. This is exactly the
	behaviour you want in a \LaTeX{} buffer.
\end{note}

\hrulebar

\section{VimTeX}\label{sec:setup-vimtex}

\texttt{vimtex} \cite{vimtex} turns Neovim into a \LaTeX{} IDE: background compilation, forward
and inverse search with the PDF viewer, a document outline, \LaTeX-aware motions
and text objects, and error parsing into the quickfix list.

The configuration in use here, from \file{lua/plugins/latex.lua}:

\begin{luacode}
{
  "lervag/vimtex",
  ft = { "tex" },                          -- load only for .tex buffers
  init = function()
    vim.g.vimtex_view_method = "zathura"   -- forward search into Zathura
    vim.g.vimtex_compiler_progname = "nvr" -- inverse search back into Neovim
    vim.g.vimtex_compiler_method = "latexmk"
    vim.g.vimtex_compiler_latexmk = {
      backend = "nvim",
      build_dir = "build",                 -- keep aux files out of the repo
      callback = 1,
      continuous = 1,                      -- recompile on every write
      executable = "latexmk",
      options = {
        "-xelatex",                        -- required: this template needs XeTeX
        "-shell-escape",
        "-interaction=nonstopmode",
        "-synctex=1",                      -- enables forward/inverse search
      },
    },
  end,
}
\end{luacode}

\begin{note}
	VimTeX's options must be set \emph{before} the plugin loads, which is why they
	go in \texttt{init} rather than \texttt{config} or \texttt{opts} --- setting
	\texttt{vim.g.vimtex\_*} after the fact silently does nothing. This is the
	single most common VimTeX configuration mistake.
\end{note}

\subsection{Commands}

VimTeX's prefix is \key{<localleader>} \key{l}, and since LazyVim leaves
\cmd{maplocalleader} at its default, that is \key{\textbackslash} \key{l}.

\begin{keys}[0.26]
	\key{\textbackslash} \key{l} \key{l} & Start / stop continuous compilation \\
	\key{\textbackslash} \key{l} \key{v} & Forward search: jump Zathura to the cursor's line \\
	\key{\textbackslash} \key{l} \key{k} & Stop the compiler                   \\
	\key{\textbackslash} \key{l} \key{e} & Show compilation errors             \\
	\key{\textbackslash} \key{l} \key{q} & Open the raw log                    \\
	\key{\textbackslash} \key{l} \key{c} & Clean auxiliary files               \\
	\key{\textbackslash} \key{l} \key{t} & Table of contents                   \\
	\key{\textbackslash} \key{l} \key{i} & Info about the current document     \\
	\key{\textbackslash} \key{l} \key{s} & Toggle which file counts as the main one \\
\end{keys}

With \texttt{continuous = 1}, \key{\textbackslash} \key{l} \key{l} once per
session is enough: every \cmd{:w} triggers a rebuild and Zathura reloads. In the
other direction, \key{<C-click>} in Zathura jumps Neovim to the corresponding
source line --- that is what \texttt{nvr} is for.

\subsection{\LaTeX-aware editing}

These are the reason to use VimTeX even if you compile in a separate terminal:

\begin{keys}[0.26]
	\key{i} \key{e} / \key{a} \key{e} & Inside / around an \texttt{environment} \\
	\key{i} \key{\$} / \key{a} \key{\$} & Inside / around inline math          \\
	\key{i} \key{c} / \key{a} \key{c} & Inside / around a command              \\
	\key{i} \key{d} / \key{a} \key{d} & Inside / around a delimiter pair       \\
	\key{i} \key{m} / \key{a} \key{m} & Inside / around a math group           \\
	\key{c} \key{s} \key{e}           & Change the surrounding environment's name \\
	\key{d} \key{s} \key{e}           & Delete the surrounding environment, keep the body \\
	\key{c} \key{s} \key{c} / \key{d} \key{s} \key{c} & Change / delete the surrounding command \\
	\key{c} \key{s} \key{\$} / \key{d} \key{s} \key{\$} & Change / delete surrounding math \\
	\key{c} \key{s} \key{d} / \key{d} \key{s} \key{d} & Change / delete surrounding delimiters \\
	\key{t} \key{s} \key{s}           & Toggle the starred form of an environment \\
	\key{t} \key{s} \key{f}           & Toggle between \cmd{\textbackslash frac} and \cmd{/} \\
	\key{t} \key{s} \key{d}           & Toggle \cmd{\textbackslash left\textbackslash right} on a delimiter pair \\
	\key{]} \key{]} in Insert mode    & Close the innermost open environment or delimiter \\
	\key{\%}                          & Extended to match \cmd{\textbackslash begin}/\cmd{\textbackslash end}, \cmd{\textbackslash if}/\cmd{\textbackslash fi} \\
\end{keys}

\begin{eg}
	The cursor is somewhere inside a \texttt{theorem} environment that should have
	been a \texttt{lemma}. Press \key{c} \key{s} \key{e}, type \texttt{lemma},
	\key{<CR>} --- both the \cmd{\textbackslash begin} and the
	\cmd{\textbackslash end} are rewritten. To delete a stray
	\cmd{\textbackslash textbf\{\dots\}} but keep the text inside it, \key{d}
	\key{s} \key{c}. To wrap the visual selection in a new environment, \key{<F6>}.
\end{eg}

\section{LuaSnip}\label{sec:setup-luasnip}

\texttt{LuaSnip} \cite{luasnip} expands abbreviations into templates with tab
stops. The
snippets here live in \file{lua/snippets/tex.lua} and are loaded by

\begin{luacode}
require("luasnip.loaders.from_lua").load({ paths = "~/.config/nvim/lua/snippets" })
\end{luacode}

A snippet is a table of \vocab{nodes}: literal text, numbered insert points, and
the formatting that stitches them together.

\begin{luacode}
local ls = require("luasnip")
local s, t, i = ls.snippet, ls.text_node, ls.insert_node
local fmta = require("luasnip.extras.fmt").fmta

return {
  -- type "dm" -> a display-math block, cursor at <1>
  s({ trig = "dm", name = "display Math" }, fmta([[
\[
    <content>
\]<cursor>
]], { content = i(1), cursor = i(2, "$0") })),

  -- type "rmk" -> a remark environment
  s({ trig = "rmk" }, fmta([[
\begin{<remark>}
    <content>
\end{<remark>}
]], { remark = t("remark"), content = i(1) })),
}
\end{luacode}

\texttt{fmta} takes a template string in which \texttt{<name>} marks a hole ---
angle brackets rather than the usual braces, precisely so that \LaTeX's own
braces need no escaping. \texttt{i(1)}, \texttt{i(2)}, \dots{} are the tab stops,
visited in order; \key{<Tab>} moves forward and \key{<S-Tab>} back.

\begin{keys}[0.30]
	\key{<Tab>}       & Expand the trigger, or jump to the next tab stop \\
	\key{<S-Tab>}     & Jump to the previous tab stop                    \\
	\key{<C-e>}       & Cycle through a choice node                      \\
	\cmd{:LuaSnipListAvailable} & List what is loaded for this filetype  \\
\end{keys}

\begin{note}[Reloading while you write snippets]
	\texttt{from\_lua} caches. After editing the snippet file, re-run the loader
\begin{vimcode}
:lua require("luasnip.loaders.from_lua").load({
  paths = "~/.config/nvim/lua/snippets" })
\end{vimcode}
	or just restart --- otherwise you will be testing the old version and
	concluding your snippet is broken.
\end{note}

\begin{remark}
	Three separate files in this configuration call a LuaSnip loader:
	\begin{itemize}[nosep, leftmargin=*]
		\item \file{init.lua} --- absolute path, works;
		\item \file{lua/plugins/luasnip-snippets.lua} --- absolute path, works;
		\item \file{lua/plugins/latex.lua} --- \emph{relative} path,
		      \file{./lua/snippets}.
	\end{itemize}
	The last one resolves against the current working directory rather than the
	config directory, so it finds nothing unless Neovim was started from
	\file{\textasciitilde/.config/nvim}. Harmless, since the other two use the
	absolute path, but worth collapsing into a single loader call the next time you
	touch the file.
\end{remark}

\section{The \texorpdfstring{\LaTeX{}}{LaTeX} Workflow, End to End}\label{sec:setup-workflow}

\begin{enumerate}[(1)]
	\item \cmd{nvim Vim.tex} --- VimTeX loads on the \texttt{tex} filetype.
	\item \key{\textbackslash} \key{l} \key{l} starts \texttt{latexmk} in continuous
	      mode; the PDF appears in Zathura.
	\item Write. Every \cmd{:w} (or \key{<C-s>}) recompiles in the background;
	      errors land in the quickfix list, reachable with \key{]} \key{q}.
	\item \key{\textbackslash} \key{l} \key{v} jumps the viewer to the cursor;
	      \key{<C-click>} in the viewer jumps the cursor back.
	\item \key{\textbackslash} \key{l} \key{t} opens the table of contents to
	      navigate between sections.
	\item \key{\textbackslash} \key{l} \key{c} cleans up before committing.
\end{enumerate}

\begin{note}
	This template requires XeTeX or LuaTeX --- it loads \texttt{fontspec} and
	\texttt{xeCJK} for the CJK fonts. That is why \texttt{vimtex\_compiler\_latexmk}
	passes \cmd{-xelatex}; with the default \texttt{pdflatex} the header aborts
	immediately with an explicit error.
\end{note}

\begin{moral}
	A configuration is a tool you maintain, not a possession you collect. Read
	\cmd{:checkhealth} when something breaks, \cmd{:verbose} when something is
	mysterious, and delete anything you have not used in a month.
\end{moral}
